\ReportSection{1. Summary of Deliverables}
\ReportSubsection{1.1 Project Objective and Scope}
We tested Todo Manager 1.5.5 using stories about how a student organizes coursework
and personal tasks. The API stores todos, projects, categories, and the links
between them. In Part A, we explored these capabilities individually. In Part B,
we used them together to test goals such as recording a task, adding it to a
course project, and checking what work remains.

\begin{ReportTable}{@{}P{.28\linewidth}P{.12\linewidth}Y@{}}
\toprule
\textbf{Measure} & \textbf{Count} & \textbf{Explanation}\\
\midrule
User stories & 20 & Goals covering todos, categories, projects, and their relationships.\\
\textit{Scenario Outlines} & 72 & Normal, alternate, and error flows.\\
Executed examples & 146 & One case for each row in the \textit{Examples} tables.\\
Passed examples & 122 & Cases that met their acceptance criteria.\\
Failed examples & 24 & Cases linked to nine findings in the bug report.\\
Policy-related failures & 16 & A subset of the 24 failures, checking workflow rules the API does not clearly promise.\\
\bottomrule
\end{ReportTable}

\ReportSubsection{1.2 Deliverables}
\begin{ReportTable}{@{}P{.24\linewidth}Y@{}}
\toprule
\textbf{Deliverable} & \textbf{Location and role}\\
\midrule
Written report & Describes the story coverage, test suite, reproduction steps, execution results, and findings.\\
Story test suite & \Path{ECSE429-Project-B/}: feature files, step definitions, shared helpers, setup code, dependency files, and the application JAR.\\
Bug report & \Path{BUG_REPORT.md}: expected and actual behavior, potential impact, reproduction steps, and related story IDs.\\
Execution records & \Path{Verification/}: Cucumber results, console logs, source hashes, and additional API checks. Section~8 lists the files used in this report.\\
\bottomrule
\end{ReportTable}

\ReportSubsection{1.3 Connecting the Tests and Findings}
\begin{center}
\begin{tikzpicture}[node distance=10bp]
\node[process,text width=.195\linewidth] (story)
{\textbf{Story N}\\Feature and steps};
\node[process,text width=.195\linewidth,right=of story] (example)
{\textbf{Outline + row}\\One test case};
\node[process,text width=.195\linewidth,right=of example] (result)
{\textbf{JSON result}\\Outcome + cleanup};
\node[process,text width=.195\linewidth,right=of result] (finding)
{\textbf{Bug ID}\\Finding and steps};
\draw[arrow] (story)--(example);
\draw[arrow] (example)--(result);
\draw[arrow] (result)--node[above,font=\fontsize{7bp}{8bp}\selectfont]{fail}(finding);
\end{tikzpicture}
\end{center}
Source hashes identify the code used for the run. Bug IDs connect failed cases
to their explanations in \Path{BUG_REPORT.md}.

\ReportSubsection{1.4 Team Participation}
\begin{ReportTable}{@{}P{.25\linewidth}P{.10\linewidth}P{.20\linewidth}Y@{}}
\toprule
\textbf{Member} & \textbf{Stories} & \textbf{Area} & \textbf{Other responsibilities}\\
\midrule
James Helou & 1--5 & Todos & Harness and shared steps\\
Jeffrey Lim & 6--10 & Categories & Written report\\
Mahan Maalekian & 11--15 & Interoperability & Integration and submission ZIP\\
Matthew Petruzziello & 16--20 & Projects & Videos\\
\bottomrule
\end{ReportTable}
